\section{Dimensions and locality of the original source network}
\label{sec:peps_source_network_dimensions}

Let $\mathcal W:a\to b$ be a chronological source circuit on an arbitrary
party set $P$. Write $G$ for its original gate occurrences with nonempty
participant sets $C_g$, and $\Xi_g$ for their finite branch sets. Let
$r_g\in C_g$ be the root chosen from the recorded participant set.
At each occurrence, the star joins $r_g$ to every other member of $C_g$
and shares the complete ordered ket and bra label in $\Xi_g\times\Xi_g$.
Every unordered pair of distinct members of $C_g$ also has its original
source sample link, with sample alphabet $[k]=\{0,\ldots,k-1\}$.
Links retain their original occurrence and kind, including parallel links.
Write $g(e)$ for the occurrence owning $e$, $D_e$ for its virtual alphabet,
and $\partial e$ for its set of incident parties. Neither finiteness nor
nonemptiness of $P$ is assumed.

\begin{theorem}[Exact dimensions of the original density links]
    \label{thm:peps_source_density_link_card}
    \lean{TNLean.Lattice.PartyCircuit.SourceCircuit.card_alphabet_source_density}
    \leanok
    \uses{def:peps_source_gate_root,def:compression_distributed_links,
        def:compression_operator_alphabets}
    As an auxiliary cardinality identity for
    \cite[proof of Theorem~5.2]{OpenAI2026PolynomialPEPS}, for any sample
    size $k\in\N$ and any constructed link $e$, including empty branch sets,
    $k=0$, and an empty link family,
    \begin{align}
        |D_e| & =
        \begin{cases}
            |\Xi_{g(e)}|^2, & e\text{ is a star link},   \\
            k,              & e\text{ is a sample link}.
        \end{cases}
        \label{eq:peps_source_density_link_card}
    \end{align}
    % Source: 04-compression.tex, lines 565--579.
\end{theorem}

\begin{proof}
    \leanok
    \uses{thm:compression_operator_alphabet_card}
    Because a star carries $\Xi_{g(e)}\times\Xi_{g(e)}$ of cardinality
    $|\Xi_{g(e)}|^2$ and a source sample link carries $[k]$,
    substituting these cardinalities into the exact alphabet formula gives
    the stated formula.
\end{proof}

\begin{theorem}[A conditional uniform bound for the original density links]
    \label{thm:peps_source_density_link_card_bound}
    \lean{TNLean.Lattice.PartyCircuit.SourceCircuit.card_alphabet_source_density_le}
    \leanok
    \uses{def:peps_source_gate_root,def:compression_distributed_links,
        def:compression_operator_alphabets}
    For the cardinality estimate used in
    \cite[proof of Theorem~5.2]{OpenAI2026PolynomialPEPS}, let $k,M\in\N$
    and assume the explicit uniform bound $|\Xi_g|\le M$ for every original
    nonempty gate occurrence $g\in G$, without requiring positivity of
    $M,k$ or the branch counts. Then every constructed link satisfies
    $|D_e|\le\max\{M^2,k\}$.
    % Source: 04-compression.tex, lines 565--588.
\end{theorem}

\begin{proof}
    \leanok
    \uses{thm:compression_operator_alphabet_bound}
    Monotonicity of multiplication in $\N$ gives
    $|\Xi_g\times\Xi_g|=|\Xi_g|^2\le M^2$, so applying the uniform
    alphabet bound with gate alphabet bound $M^2$ and sample size $k$
    yields the inequality.
\end{proof}

\begin{theorem}[The endpoints are original participants of the owner gate]
    \label{thm:peps_source_density_link_locality}
    \lean{TNLean.Lattice.PartyCircuit.SourceCircuit.distributedLinkParties_source_subset}
    \leanok
    \uses{def:peps_source_gate_root,def:compression_distributed_links,
        def:compression_distributed_link_parties}
    For the locality of the original links in
    \cite[proof of Theorem~5.2]{OpenAI2026PolynomialPEPS}, with repeated
    occurrences and parallel links distinguished and root membership
    inherited from each nonempty participant set rather than assumed
    separately, every constructed link satisfies $\partial e\subseteq C_{g(e)}$.
    % Source: 04-compression.tex, lines 581--588.
\end{theorem}

\begin{proof}
    \leanok
    \uses{thm:peps_source_gate_root_mem,thm:compression_original_gate_locality}
    Because the chosen root satisfies $r_g\in C_g$ at every original
    occurrence, original-gate locality of the distributed link
    construction applies with these derived root memberships.
\end{proof}

\section{Incidence of the density links}
\label{sec:peps_source_density_link_incidence}

Let $\mathcal W:a\to c$ be a chronological source circuit on a party set
$P$. Write $G$ for its original nonempty gate occurrences and $C_g$ for
the participant set of occurrence $g$. The star root belongs to $C_g$.
At each occurrence the density network has its star links and its original
pair-sample links. Parallel links are retained as distinct links.

\begin{theorem}[The retained occurrences preserve participation counts]
    \label{thm:peps_nonempty_gate_participation}
    \lean{TNLean.Lattice.PartyCircuit.SourceCircuit.card_incidentNonemptyGates_eq_participationCount}
    \leanok
    \uses{def:compression_lifetime_gates}
    Tracking the original participation count in
    \cite[proof of Theorem~5.2]{OpenAI2026PolynomialPEPS} without assuming
    finiteness or nonemptiness of the party set, the number of occurrences
    $g\in G$ with $p\in C_g$ for any party $p$ equals its full
    participation count in $\mathcal W$.
    % Source: 04-compression.tex, lines 140--142 and 581--588.
\end{theorem}

\begin{proof}
    \leanok
    Removing an occurrence with empty participant set removes no occurrence
    involving $p$. Conversely, an occurrence involving $p$ has nonempty
    participant set. These two maps give a bijection preserving the original
    occurrence, and hence preserve the count.
\end{proof}

\begin{theorem}[Constant incidence of the density links]
    \label{thm:peps_actual_density_link_incidence}
    \lean{TNLean.Lattice.PartyCircuit.SourceCircuit.card_incidentSourceDensityLinks_le}
    \leanok
    \uses{def:peps_source_gate_root,def:compression_distributed_links,
        thm:peps_nonempty_gate_participation}
    For the incidence estimate in
    \cite[proof of Theorem~5.2]{OpenAI2026PolynomialPEPS}, suppose $b\in\N$,
    every occurrence $g\in G$ has $|C_g|\le b$, and every party has full
    participation count at most $b$. Counting star and sample links
    (including parallel links) separately without any additional
    root-membership hypothesis, the number of density-network links
    incident to any party is at most $2b(b-1)$.
    % Source: 04-compression.tex, lines 581--588.
\end{theorem}

\begin{proof}
    \leanok
    \uses{thm:peps_source_gate_root_mem,thm:peps_nonempty_gate_participation,
        thm:compression_distributed_incidence_bound}
    The chosen root lies in its original participant set. Each incident
    occurrence contributes at most $2(b-1)$ star and sample links, and at
    most $b$ occurrences involve the party. Multiplying the two bounds
    gives the result. An empty gate family is included.
\end{proof}

\section{Finite coordinates for the local memories}

The coordinate spaces in this section are the memories of the original
ordered registers. The assertions are auxiliary consequences of finite
dimensionality for the construction in
\cite[proof of Theorem~5.2]{OpenAI2026PolynomialPEPS},
\texttt{04-compression.tex}, lines 17--19, 253--267, 409--434 and 565--588.
They require no numerical bound on a private register dimension.

\begin{theorem}[Finite endpoint registers in fixed source slots]
    \label{thm:peps_finite_source_slot_registers}
    \lean{TNLean.Lattice.PartyCircuit.SourceInventory.finiteDimensional_registers_slotLayout}
    \leanok
    \uses{def:peps_source_slots}
    Let $R$ be a finite ordered list of pair-source occurrences. For each
    occurrence $i$, let $U_i$ and $V_i$ be finite-dimensional complex
    Hilbert spaces. Every register in the fixed ordered list consisting
    of the left register with space $U_i$ followed by the right register
    with space $V_i$, for each $i$ in order, is finite dimensional.
\end{theorem}

\begin{proof}
    \leanok
    \uses{def:peps_source_slots}
    We proceed by induction on the list of occurrences: the empty list has
    no registers, while removing the first occurrence removes its two
    endpoint registers with spaces $U_0$ and $V_0$ and leaves the
    remaining registers to the induction hypothesis.
\end{proof}

\begin{theorem}[Finite registers in the all-party source layout]
    \label{thm:peps_finite_all_party_source_registers}
    \lean{TNLean.Lattice.PartyCircuit.SourceCircuit.finiteDimensional_allResidualSource_registers}
    \leanok
    \uses{def:peps_all_party_source_registers,thm:peps_finite_source_slot_registers}
    Let $W$ be an original chronological source composition with arbitrary
    party set $P$, and let $R_W$ be its fixed all-party source layout.
    Every register in $R_W$ is finite dimensional. The original endpoint
    dimensions may be zero, and the source inventory may be empty.
\end{theorem}

\begin{proof}
    \leanok
    \uses{def:peps_all_party_source_registers,thm:peps_finite_source_slot_registers}
    Because each endpoint space is an original Euclidean space
    $\C^{u_e}$ or $\C^{v_e}$, applying the preceding theorem in the
    original order of source occurrences establishes finite dimensionality.
\end{proof}

\begin{definition}[Finite orthonormal coordinates of a local source memory]
    \label{def:peps_all_party_local_source_basis}
    \lean{TNLean.Lattice.PartyCircuit.SourceCircuit.allResidualSourceBasis}
    \leanok
    \uses{thm:peps_finite_all_party_source_registers,
        thm:peps_finite_register_memories,def:peps_party_registers}
    With $W$ and $R_W$ as above, write $Q=P\sqcup\{*\}$ for the grouped
    owner set. For $q\in Q$, let $R_W(q)$ be the ordered sublist of source
    registers owned by $q$, and let $H_q=\mathcal H_{R_W(q)}$ be its
    tensor-product memory. Choose an orthonormal basis
    \begin{align}
        (b_{q,i})_{0\le i<\dim_{\C}H_q}
    \end{align}
    of $H_q$. This is a finite basis of the selected memory.
\end{definition}

\begin{definition}[Finite orthonormal coordinates of a local discarded memory]
    \label{def:peps_all_party_local_discard_basis}
    \lean{TNLean.Lattice.PartyCircuit.SourceCircuit.allPartyDiscardBasis}
    \leanok
    \uses{thm:peps_finite_register_memories,def:peps_affected_owner,
        def:peps_party_registers}
    Let $G$ be the original finite ordered list of discarded registers,
    and suppose that the space of every register in $G$ is finite
    dimensional. Relabel each original owner $p$ by its copy in
    $Q=P\sqcup\{*\}$. For $q\in Q$, let $G(q)$ be the ordered sublist
    of the relabelled registers owned by $q$, and put
    $K_q=\mathcal H_{G(q)}$. Choose an orthonormal basis
    \begin{align}
        (c_{q,j})_{0\le j<\dim_{\C}K_q}
    \end{align}
    of $K_q$. Relabelling leaves every register space unchanged, and
    selection preserves finite dimensionality of the tensor-product
    memory, so this basis requires only the stated per-register
    finite-dimensionality assumption.
\end{definition}

A zero-dimensional memory has an empty basis. By contrast, the memory
of an empty register list is $\C$, and its basis has one element.
The dimensions appearing in these basis index sets are not estimates
on the private memory dimensions.

\begin{theorem}[Distributed physical approximation from local gate expansions]
    \label{thm:peps_approximate_sampled_distributed_error}
    \lean{TNLean.Lattice.PartyCircuit.EffectCircuit.exists_sampled_contractedOperator_approximate_error_le}
    \leanok
    \uses{def:peps_approximate_original_operators,
        thm:peps_approximate_source_sample,
        thm:peps_actual_sampled_distributed_contraction,
        def:peps_all_party_local_source_basis,def:peps_all_party_local_discard_basis}
    Let $P$ be a finite party set, and list its members without repetition.
    For each $p\in P$, let $H_p$ be its finite-dimensional whole initial
    memory and let $x_p\in H_p$ satisfy $\|x_p\|=1$. Put
    $\psi=\bigotimes_{p\in P}x_p$. Retain the original chronological
    composition and its gate operators $G_g$. Its output has the original
    physical register of dimension $d_p$ at each party, followed by a
    finite ordered list of discarded private registers, each with a
    finite-dimensional space.

    Suppose that the original number of nonprivate occurrences is at
    most $M$ and that each party participates in at most $\ell$ such
    occurrences. Let $B,D\geq1$, with $d_p\leq D$ for every party.
    Fix $0<\varepsilon\leq1$ and put
    \begin{align}
        \delta_a=\frac{\varepsilon}{16\max\{1,M\}}.
    \end{align}
    Suppose the monomial list at each gate is given at this
    accuracy: its monomials are allowed, each contains at most $r$ pair
    effects, their absolute coefficient sum is at most $B$, and their
    sum differs from $G_g$ in operator norm by at most $\delta_a$.
    Assume that the original private operations and the original $G_g$
    are contractions. Retain the prescribed participating sets, whose
    cardinalities are different from one; the empty set is permitted.
    These are the local approximation conditions of
    Definition~\ref{def:peps_approximate_original_operators}.

    Rescale every monomial list by $(1+\delta_a)^{-1}$ to obtain the
    original contraction circuit $W$. Set
    \begin{align}
        \delta_e=\frac{\varepsilon}{8\max\{1,M\}},
        \qquad
        m=\left\lceil\left(\frac{rB}{\delta_e}\right)^2\right\rceil+1,
    \end{align}
    and form its physical-first source replacement $R$ with stack
    length $m$. If $N_R$ is the number of source occurrences in
    $R$, set
    \begin{align}
        Q=B^{4\ell}D^4,
        \qquad
        k=\left\lceil
        \left(\frac{8\max\{1,N_R\}Q}{\varepsilon}\right)^2
        \right\rceil+1.
    \end{align}
    Choose the local source and discarded-memory orthonormal bases of
    Definitions~\ref{def:peps_all_party_local_source_basis} and
    \ref{def:peps_all_party_local_discard_basis}. The discarded list is
    the auxiliary list of $W$ followed by the original private
    list; its register-wise finite dimensionality follows from the
    original construction and the given assumptions.

    There exist normalized Schmidt coefficients and two isometric
    endpoint frames for every source occurrence and branch of
    $R$, representing its original source vector in its original
    endpoint dimensions. There is one Gaussian sample array $\omega$,
    with $k$ samples per source occurrence, such that the following
    holds. Construct the incident physical matrices, derived gate roots,
    and distributed contraction $\mathcal C_R(\omega)$ of
    Theorem~\ref{thm:peps_actual_sampled_distributed_contraction}. Its gate
    labels are ordered ket and bra branch pairs, and the coefficient at
    a nonempty occurrence $g$ is
    \begin{align}
        c_g(\xi,\zeta)
        =k^{-\binom{|C_g|}{2}}
        a_g(\xi)\overline{a_g(\zeta)}.
    \end{align}
    Every source sample assignment is common to its two endpoints.
    Let $\lambda_R$ be the scalar contributed by the empty-participant
    gate occurrences and put $\gamma_R=\lambda_R\overline{\lambda_R}$.
    For the given local lists at the prescribed accuracy in
    \cite[proof of Theorem~5.2]{OpenAI2026PolynomialPEPS}
    (\texttt{04-compression.tex}, lines 17--19, 137--175, 199--229,
    342--381, 565--588 and 590--600), with no numerical bound on the
    original private dimensions and with the scalar remaining
    multiplicative even when it vanishes,
    \begin{align}
        \bigl\|\gamma_R\mathcal C_R(\omega)-\rho_G\bigr\|_1
        \leq\varepsilon,
    \end{align}
    where $\rho_G$ is the physical partial trace of the original
    chronological output on $\psi$, evaluated with the original
    operators $G_g$ and expressed in the original party-labelled
    physical coordinates.
\end{theorem}

\begin{proof}
    \leanok
    \uses{thm:peps_approximate_source_sample,
        thm:peps_actual_sampled_distributed_contraction,
        thm:peps_finite_whole_party_input,thm:peps_auxiliary_register_finiteness,
        def:peps_all_party_local_source_basis,def:peps_all_party_local_discard_basis}
    Finite dimensionality of the whole initial party memories yields
    finite orthonormal input coordinates, allowing
    Theorem~\ref{thm:peps_approximate_source_sample} to be applied with
    the recorded local approximations. That theorem provides the
    normalized Schmidt data and the sample array for which the sampled
    physical operator differs from $\rho_G$ in trace norm by at most
    $\varepsilon$. In the stated finite local bases, the exact
    contraction identity of
    Theorem~\ref{thm:peps_actual_sampled_distributed_contraction}
    identifies this same sampled operator with
    $\gamma_R\mathcal C_R(\omega)$, and substituting this equality into
    the physical error estimate completes the proof.
\end{proof}

\section{Original resources of the sampled contraction}

The network in this section is the same network used in the sampled
physical-density identity: its gates are the original nonempty occurrences,
its roots are chosen from their recorded participant sets, and its links
retain their original occurrence and kind. No root on an empty participant
set is required.

\begin{theorem}[Real power bounds for the density alphabets]
    \label{thm:peps_actual_density_alphabet_power}
    \lean{TNLean.Lattice.PartyCircuit.SourceCircuit.card_alphabet_source_density_le_of_power_bounds}
    \leanok
    \uses{thm:peps_source_density_link_card,def:peps_source_gate_root}
    Let $R$ be a chronological source composition on an arbitrary party
    set. Write $\Xi_g$ for the branch labels of its nonempty gate occurrence
    $g$, and use the recorded roots and pair-sample links. Suppose
    $L\geq1$, $C_M,C_k\geq0$, $u,v\in\N$, and $k\in\N$ satisfy
    $|\Xi_g|\leq C_M L^u$ for every such $g$ and $k\leq C_k L^v$.
    Then, as used in
    \cite[proof of Theorem~5.2]{OpenAI2026PolynomialPEPS}
    (\texttt{04-compression.tex}, lines 565--588) and including empty
    branch alphabets, $k=0$, and an empty gate family, every alphabet
    $D_e$ of this density network satisfies
    $|D_e|\leq(C_M^2+C_k)L^{\max\{2u,v\}}$.
\end{theorem}

\begin{proof}
    \leanok
    \uses{thm:peps_source_density_link_card}
    On a star, $|D_e|=|\Xi_{g(e)}|^2\leq C_M^2 L^{2u}$.
    On a sample link, $|D_e|=k\leq C_k L^v$.
    Since $L\geq1$, each is bounded by
    $(C_M^2+C_k)L^{\max\{2u,v\}}$.
\end{proof}

\begin{theorem}[Original polynomial bounds for the contraction alphabets]
    \label{thm:peps_actual_physicalFirst_density_polynomial}
    \lean{TNLean.Lattice.PartyCircuit.OriginalCircuit.card_physicalFirstDensityAlphabet_le_power}
    \leanok
    \uses{thm:peps_actual_density_alphabet_power,
        thm:peps_physical_first_branch_polynomial,
        thm:peps_physical_first_sample_polynomial}
    Let $W:a\to H_{\mathrm{phys}}\otimes H_{\mathrm{priv}}$ be an original
    chronological circuit. Suppose its monomial lists have at most $K$
    terms, at most $r$ pair effects per term, and absolute coefficient
    sum at most $B$. Suppose its nonprivate gate count is at most $N$,
    and every original participant set has cardinality at most $b$.
    Let $L\geq1$, $B,D\geq0$, $C_K,C_B\geq0$, $C_N\geq1$, and
    $C_D\in\R$ satisfy
    $K\leq C_KL^u$, $N\leq C_NL^n$, $B\leq C_BL^s$ and
    $D\leq C_DL^d$, where $u,n,s,d,p\in\N$.
    Put $\delta=(L^p)^{-1}/(8\max\{1,N\})$ and suppose $\delta>0$.
    Form the physical-first replacement $R$ at tolerance $\delta$.
    For any natural parameter $\ell$, use the sample count prescribed
    by $Q=B^{4\ell}D^4$ and accuracy $(L^p)^{-1}$.
    Define
    \begin{align}
        C_s & =\max\{1,C_N\tbinom b2\},                  &
        C_b & =C_K\bigl(64C_N^2(rC_B)^2+2\bigr)^r,\notag   \\
        C_k & =64C_s^2(C_B^{4\ell}C_D^4)^2+2,            &
        E   & =\max\{2(u+2r(n+s+p)),2(n+4\ell s+4d+p)\}.
        \label{eq:peps_actual_density_power_constants}
    \end{align}
    Then, for the auxiliary resource estimate in
    \cite[proof of Theorem~5.2]{OpenAI2026PolynomialPEPS}
    (\texttt{04-compression.tex}, lines 137--175, 199--229 and 565--588),
    without any numerical bound on a private dimension or a source
    Schmidt rank, every alphabet of the density network on the nonempty
    occurrences of $R$ satisfies $|D_e|\leq(C_b^2+C_k)L^E$.
\end{theorem}

\begin{proof}
    \leanok
    \uses{thm:peps_actual_density_alphabet_power,
        thm:peps_physical_first_branch_polynomial,
        thm:peps_physical_first_sample_polynomial}
    Applying the proved original branch-count bound at the ambient
    occurrence of each retained nonempty gate and the proved sample
    bound to this same replacement and sample count gives
    $|\Xi_g|\leq C_b L^{u+2r(n+s+p)}$ and
    $k\leq C_k L^{2(n+4\ell s+4d+p)}$, from which the preceding theorem
    yields the stated bound.
\end{proof}

\begin{theorem}[Original incidence bound for the contraction]
    \label{thm:peps_actual_physicalFirst_density_incidence}
    \lean{TNLean.Lattice.PartyCircuit.OriginalCircuit.card_incidentPhysicalFirstDensityLinks_le}
    \leanok
    \uses{thm:peps_actual_density_link_incidence,
        def:peps_physical_first_original_locality,thm:peps_physical_first_original_locality,
        thm:peps_physical_first_occurrences,thm:peps_source_only_resources}
    Let $W$ be an original chronological circuit with expansion bound
    $(r,S)$ and let $\delta>0$. Form its physical-first source
    replacement $R$ at tolerance $\delta$. Suppose every original gate has
    at most $b$ participants and every party's original whole-lifetime
    participation count is at most $b$. As an auxiliary consequence of
    the original incidence bound in
    \cite[proof of Theorem~5.2]{OpenAI2026PolynomialPEPS}
    (\texttt{04-compression.tex}, lines 140--142 and 581--588), even when
    the party set is infinite or empty, at most $2b(b-1)$ links of the
    density network of $R$ are incident to any party.
\end{theorem}

\begin{proof}
    \leanok
    \uses{thm:peps_actual_density_link_incidence,
        def:peps_physical_first_original_locality,thm:peps_physical_first_original_locality,
        thm:peps_physical_first_occurrences,thm:peps_source_only_resources}
    The original occurrence correspondence preserves participant sets and
    hence the arity bound, while the physical-first output exchange
    preserves participation counts so that the source replacement inherits
    the original participation bound. Applying the incidence estimate for
    the nonempty-occurrence network, whose roots already belong to their
    recorded participant sets, completes the proof.
\end{proof}

\begin{theorem}[A common original gate for the contraction endpoints]
    \label{thm:peps_actual_physicalFirst_density_locality}
    \lean{TNLean.Lattice.PartyCircuit.OriginalCircuit.exists_original_gate_of_physicalFirst_density_link}
    \leanok
    \uses{thm:peps_source_density_link_locality,
        def:peps_physical_first_original_locality,thm:peps_physical_first_original_locality}
    With $W,R,\delta$ and the expansion bound as above, take any link of
    the density network of $R$ and two parties incident to it.
    Then, as used in
    \cite[proof of Theorem~5.2]{OpenAI2026PolynomialPEPS}
    (\texttt{04-compression.tex}, lines 581--588) without requiring
    finiteness or nonemptiness of the party set, some original nonprivate
    occurrence of $W$ contains both parties in its participant set.
\end{theorem}

\begin{proof}
    \leanok
    \uses{thm:peps_source_density_link_locality,
        def:peps_physical_first_original_locality,thm:peps_physical_first_original_locality}
    Because both endpoints belong to the retained link's original owner
    gate in $R$ and the original occurrence correspondence preserves the
    entire participant set, the inverse of that ambient occurrence in $W$
    contains both parties.
\end{proof}

\begin{definition}[Termwise bounds for prescribed approximate gate expansions]
    \label{def:peps_effect_termwise_bounds}
    \lean{TNLean.Lattice.PartyCircuit.EffectCircuit.IsTermwiseBounded}
    \leanok
    For each nonprivate gate occurrence, retain its prescribed ordered
    monomial expansion. Termwise bounds $(r,C)$ mean that every listed
    monomial contains at most $r$ pair effects and has scalar coefficient
    of modulus at most $C$. These conditions do not assert a bound on the
    sum of the coefficient moduli. The contraction conditions and the
    approximation to the original gate are specified separately.
\end{definition}

\begin{theorem}[A polynomial coefficient majorant for approximate gate lists]
    \label{thm:peps_effect_coefficient_majorant}
    \lean{TNLean.Lattice.PartyCircuit.EffectCircuit.isExpansionBounded_of_isTermwiseBounded}
    \lean{TNLean.Lattice.PartyCircuit.EffectCircuit.exists_one_le_isExpansionBounded_of_termwise_power_bounds}
    \leanok
    \uses{def:peps_effect_termwise_bounds,thm:peps_gate_coefficient_sum}
    Suppose each prescribed expansion at an original nonprivate gate
    occurrence has at most $K$ terms and has termwise bounds $(r,C)$,
    where $C\geq0$. Its absolute coefficient sum is at most $KC$.
    In particular, let $L\geq1$, let $C_K,C_C\geq0$, and let $u,v$ be
    nonnegative integers. If $K\leq C_KL^u$ and the termwise bounds are
    $(r,C_CL^v)$, then all these sums admit the common bound
    $B=(1+C_KC_C)L^{u+v}\geq1$.
    This bound is fixed by the original polynomial constants and
    exponents. It is independent of private dimensions and of the
    subsequent effect-elimination expansion. Empty lists are permitted.
\end{theorem}
\begin{proof}
    \leanok
    \uses{thm:peps_gate_coefficient_sum}
    Summing the individual coefficient inequalities over each original
    ordered list, multiplying its length bound by the nonnegative
    individual bound, and combining the powers of $L$ produces the
    displayed majorant, which is at least one because $L\geq1$;
    composition and the addition of spectator registers preserve the same
    local bounds.
\end{proof}

\begin{theorem}[Absorbing a scalar into one distributed local matrix]
    \label{thm:peps_distributed_operator_scalar_absorption}
    \lean{TNLean.PEPS.Approximation.DistributedOperatorContraction.contractedOperator_smul_at_party}
    \leanok
    \uses{thm:compression_operator_matrix_sum,thm:peps_party_product_scalar}
    Let $P$ be a finite party set and let $\mathcal C(T)$ be the
    distributed star-and-sample contraction with a given finite
    gate collection, participant sets, participating roots, finite
    gate alphabets, and sample size $k\in\N$. Let the local rectangular
    matrices $T_v(\xi,j)$ have arbitrary fixed row and column index sets.
    For a given party $v_0\in P$ and $\mu\in\C$, put
    $T'_v(\xi,j)=\mu T_v(\xi,j)$ when $v=v_0$ and
    $T'_v(\xi,j)=T_v(\xi,j)$ otherwise.
    Then, as used in
    \cite[proof of Theorem~5.2]{OpenAI2026PolynomialPEPS}
    (\texttt{04-compression.tex}, lines 573--579) and including $\mu=0$,
    empty gate alphabets, and $k=0$, the link family and its alphabets
    are unchanged and $\mathcal C(T')=\mu\mathcal C(T)$.
\end{theorem}

\begin{proof}
    \leanok
    \uses{thm:compression_operator_matrix_sum,thm:peps_party_product_scalar}
    Expanding the contraction as the sum over common gate labels and
    sample indices of the weighted product of local matrices, scaling the
    matrix at $v_0$ multiplies every summand by $\mu$. Distributing this
    scalar through both finite sums and commuting it with each gate
    coefficient yields the identity without division.
\end{proof}
